\documentclass[11pt,a4paper]{article}
\usepackage[margin=20mm,headheight=15pt]{geometry}
\usepackage{fontspec}
\setmainfont{TeX Gyre Pagella}
\setmonofont{DejaVu Sans Mono}[Scale=0.82]
\usepackage{amsmath}
\usepackage{unicode-math}
\setmathfont{TeX Gyre Pagella Math}
\setcounter{MaxMatrixCols}{20}
\usepackage{microtype}
\usepackage{enumitem}
\usepackage{fvextra}
\usepackage{booktabs}
\usepackage{longtable}
\usepackage{fancyhdr}
\usepackage[hidelinks]{hyperref}
\setlength{\parindent}{0pt}
\setlength{\parskip}{6pt}
\setlength{\emergencystretch}{3em}
\hyphenpenalty=10000
\exhyphenpenalty=10000
\raggedright
\DefineVerbatimEnvironment{Code}{Verbatim}{fontsize=\small,
  breaklines=true,breakanywhere=true,breaksymbolleft={},breaksymbolright={},
  frame=single,framesep=5pt}
\setlist[enumerate]{leftmargin=7mm,itemsep=3pt,topsep=5pt,parsep=0pt}
\pagestyle{fancy}
\fancyhf{}
\fancyhead[L]{\small Computer Vision 1}
\fancyhead[R]{\small Exam 3 · Questions}
\fancyfoot[C]{\small \thepage}
\hypersetup{pdftitle={Exam 3 — Synthesis and transfer},pdfsubject={CV1 practice material}}
\begin{document}
{\LARGE\bfseries Exam 3 — Synthesis and transfer}\par
{\large Questions only}\par\bigskip
\textbf{Name:} \rule{70mm}{0.4pt}\hfill \textbf{Date:} \rule{30mm}{0.4pt}
\section*{Conditions}
\begin{itemize}[leftmargin=5mm,itemsep=5pt]
\item 50 questions: 30 single-answer multiple choice and 20 numerical. Each is worth 1 raw practice point.
\item Suggested training limit: \textbf{150 minutes}. This is not a claim about the actual exam duration or question count.
\item Use up to three double-sided A4 sheets,\allowbreak{} prepared by you,\allowbreak{} and a local Jupyter notebook. No internet,\allowbreak{} AI assistance or answer key during the attempt.
\item Use Python/\allowbreak{}NumPy for computation; SciPy is also allowed. You do not need executable OpenCV or PyTorch. Specialized routines are provided when needed.
\item For multiple choice,\allowbreak{} select exactly one of A–D. Attempt every question.
\item Unless an exact integer is requested,\allowbreak{} round the final numerical answer to \textbf{two decimal places}. Keep full precision in intermediate steps. Use fractions rather than percentages unless specified.
\item Numerical practice grading uses the syllabus rule: within ±0.01 of the unrounded reference,\allowbreak{} inclusive. Counts require the exact integer; no partial credit.
\item Conventions stated in a question override defaults. Arrays use [row,\allowbreak{}column] =\allowbreak{} [y,\allowbreak{}x]; geometric vectors are columns; image-plane projection is virtual/\allowbreak{}front-facing; foreground is 1.
\item This is original,\allowbreak{} challenging training material based on the supplied course files,\allowbreak{} not a lecturer-issued mock or a prediction of the exam.
\end{itemize}
Solutions are in a separate PDF. Do not open them until you finish.
\clearpage\section*{Part A — Single-answer multiple choice}
\begin{minipage}{\linewidth}
\textbf{Q01 — Multiple choice}\par\smallskip
A unit normal \(n_c\) is expressed in the camera frame. The camera uses \(X_c=R X_w+t\). Which expression converts the normal to the world frame?
\begin{enumerate}[label=\textbf{\Alph*.}]
\item \(n_w=R^{\mathsf T}(n_c-t)\)
\item \(n_w=Kn_c\),\allowbreak{} followed by perspective division
\item \(n_w=R^{\mathsf T}n_c\); translation does not act on a direction vector.
\item \(n_w=Rn_c+t\)
\end{enumerate}
\end{minipage}
\par\vspace{12pt}
\begin{minipage}{\linewidth}
\textbf{Q02 — Multiple choice}\par\smallskip
You resize an image to one quarter of each spatial dimension by keeping every fourth pixel. Why should a suitable low-pass filter normally be applied BEFORE subsampling?
\begin{enumerate}[label=\textbf{\Alph*.}]
\item Keeping every fourth pixel is itself an ideal low-pass filter,\allowbreak{} so preprocessing cannot influence aliasing.
\item \(I_t\) guarantees exact recovery of every original sample,\allowbreak{} even when the image has frequencies above the coarse-grid limit.
\item High spatial frequencies can otherwise alias into false lower-frequency patterns.
\item Low-pass filtering after subsampling removes aliasing just as reliably as filtering before subsampling.
\end{enumerate}
\end{minipage}
\par\vspace{12pt}
\begin{minipage}{\linewidth}
\textbf{Q03 — Multiple choice}\par\smallskip
A single image is modelled as \(I=A\odot S\). Without further constraints,\allowbreak{} can it uniquely determine both albedo and shading?
\begin{enumerate}[label=\textbf{\Alph*.}]
\item No,\allowbreak{} but dividing every pixel by 255 always resolves the ambiguity.
\item No; multiple albedo/\allowbreak{}shading combinations can produce the same image.
\item Yes,\allowbreak{} because one RGB pixel always supplies all geometry and material parameters.
\item Yes,\allowbreak{} if the image has a high enough resolution,\allowbreak{} regardless of assumptions.
\end{enumerate}
\end{minipage}
\par\vspace{12pt}
\begin{minipage}{\linewidth}
\textbf{Q04 — Multiple choice}\par\smallskip
Why can directly fitting the linear photometric-stereo model to strongly gamma-compressed pixel values be wrong?
\begin{enumerate}[label=\textbf{\Alph*.}]
\item Normalizing the fitted g removes all gamma effects,\allowbreak{} so only albedo can be affected by a nonlinear response.
\item The model assumes a linear relation between measured intensity and radiance; the response must be linearized or modelled.
\item Gamma can always be absorbed into one constant camera gain,\allowbreak{} independent of intensity and light direction.
\item No correction is needed: any common gamma exponent preserves the linear equations exactly under a change of albedo alone.
\end{enumerate}
\end{minipage}
\par\vspace{12pt}
\begin{minipage}{\linewidth}
\textbf{Q05 — Multiple choice}\par\smallskip
What determines a sensor channel's response to an illuminated coloured surface?
\begin{enumerate}[label=\textbf{\Alph*.}]
\item The interaction of source spectrum,\allowbreak{} surface spectral reflectance and that channel's spectral sensitivity.
\item Only the surface's colour name,\allowbreak{} independent of light and sensor.
\item Only the maximum channel value in a different pixel.
\item Only the camera's extrinsic translation.
\end{enumerate}
\end{minipage}
\par\vspace{12pt}
\begin{minipage}{\linewidth}
\textbf{Q06 — Multiple choice}\par\smallskip
An integrable normal-derived gradient field is reconstructed on a connected image domain with known x/\allowbreak{}y spacing and at least one nonzero slope. Which ambiguity remains without an anchor or boundary height?
\begin{enumerate}[label=\textbf{\Alph*.}]
\item No height ambiguity remains,\allowbreak{} because unit-normal normalization already fixes the absolute reference height.
\item An additive constant in height.
\item An arbitrary multiplicative height scale,\allowbreak{} even though p and q and the x/\allowbreak{}y spacing are fixed.
\item An independent additive height constant for every pixel,\allowbreak{} even while all neighbouring derivatives stay fixed.
\end{enumerate}
\end{minipage}
\par\vspace{12pt}
\begin{minipage}{\linewidth}
\textbf{Q07 — Multiple choice}\par\smallskip
Apply a normalized averaging kernel to a constant positive image with zero padding. What is correct?
\begin{enumerate}[label=\textbf{\Alph*.}]
\item All outputs become zero because averaging removes the DC component.
\item Every output must be identical,\allowbreak{} including boundaries,\allowbreak{} for any padding rule.
\item Only corners become brighter because there are fewer neighbours.
\item Interior outputs remain constant,\allowbreak{} but outputs near boundaries can be darker.
\end{enumerate}
\end{minipage}
\par\vspace{12pt}
\begin{minipage}{\linewidth}
\textbf{Q08 — Multiple choice}\par\smallskip
A closing operation uses a small structuring element. Which statement is safest?
\begin{enumerate}[label=\textbf{\Alph*.}]
\item \(I_t\) can fill sufficiently small gaps/\allowbreak{}holes relative to the element's shape,\allowbreak{} but need not fill a large cavity.
\item Closing removes every isolated foreground component without changing holes.
\item Closing means erosion followed by dilation.
\item One closing always fills every hole in every binary object.
\end{enumerate}
\end{minipage}
\par\vspace{12pt}
\begin{minipage}{\linewidth}
\textbf{Q09 — Multiple choice}\par\smallskip
Erode an all-ones finite image using a centred 3 by 3 all-ones element with outside-image values defined as zero. What happens?
\begin{enumerate}[label=\textbf{\Alph*.}]
\item Its foreground expands by one pixel outside the image.
\item \(I_t\) remains all ones because every stored pixel is one.
\item Its outermost pixel border is removed.
\item Every interior pixel becomes zero while the border stays one.
\end{enumerate}
\end{minipage}
\par\vspace{12pt}
\begin{minipage}{\linewidth}
\textbf{Q10 — Multiple choice}\par\smallskip
Circularly translate a discrete image by an integer offset. How does the DFT transform under the standard shift theorem?
\begin{enumerate}[label=\textbf{\Alph*.}]
\item Only magnitude changes; phase is fixed.
\item The spectrum is only reordered by fftshift; individual frequency coefficients never change.
\item Both magnitude and phase must stay unchanged because the image content is the same.
\item Each coefficient is multiplied by a frequency-dependent phase factor of unit magnitude; its magnitude is unchanged.
\end{enumerate}
\end{minipage}
\par\vspace{12pt}
\begin{minipage}{\linewidth}
\textbf{Q11 — Multiple choice}\par\smallskip
For a real-valued image,\allowbreak{} what relationship holds between opposite-frequency DFT coefficients?
\begin{enumerate}[label=\textbf{\Alph*.}]
\item Their imaginary components must all be zero,\allowbreak{} because the original image contains only real values.
\item They are complex conjugates,\allowbreak{} with indices interpreted modulo the image size.
\item They are always identical in both real and imaginary components,\allowbreak{} so their phases must also be identical.
\item They are always negatives of each other,\allowbreak{} so the spectrum is antisymmetric at every opposite-frequency pair.
\end{enumerate}
\end{minipage}
\par\vspace{12pt}
\begin{minipage}{\linewidth}
\textbf{Q12 — Multiple choice}\par\smallskip
A detected intensity edge does not necessarily mark a depth discontinuity. Why?
\begin{enumerate}[label=\textbf{\Alph*.}]
\item Canny hysteresis guarantees that only physical depth boundaries survive its high threshold.
\item A material boundary cannot affect intensity,\allowbreak{} so every strong edge must be a depth discontinuity.
\item Increasing Gaussian \(\sigma\) makes intensity edges uniquely attributable to changes in depth.
\item Reflectance or illumination changes can also cause an intensity edge.
\end{enumerate}
\end{minipage}
\par\vspace{12pt}
\begin{minipage}{\linewidth}
\textbf{Q13 — Multiple choice}\par\smallskip
Why does SIFT search across a Gaussian scale space and use Difference-of-Gaussians extrema?
\begin{enumerate}[label=\textbf{\Alph*.}]
\item To locate distinctive structures at their characteristic scales; DoG approximates a scale-normalized LoG response.
\item To recover absolute object size in metres from one image.
\item To guarantee invariance to arbitrary projective deformation of every patch.
\item To force all keypoints to have identical pixel scale.
\end{enumerate}
\end{minipage}
\par\vspace{12pt}
\begin{minipage}{\linewidth}
\textbf{Q14 — Multiple choice}\par\smallskip
For positive affine brightness changes \(I'=aI+b\) with \(a>0\),\allowbreak{} why can a normalized gradient-based SIFT descriptor be relatively robust?
\begin{enumerate}[label=\textbf{\Alph*.}]
\item Dividing pixel coordinates by the focal length removes b.
\item Gradients remove the constant offset b,\allowbreak{} and descriptor normalization reduces the common gain a.
\item The descriptor stores the original unnormalized RGB values exactly.
\item SIFT estimates and removes every spatially varying shadow without assumptions.
\end{enumerate}
\end{minipage}
\par\vspace{12pt}
\begin{minipage}{\linewidth}
\textbf{Q15 — Multiple choice}\par\smallskip
An affine RANSAC sample contains three correct correspondences,\allowbreak{} but the three source points are collinear. Why should the sample be rejected?
\begin{enumerate}[label=\textbf{\Alph*.}]
\item Three pairs supply six scalar equations,\allowbreak{} so the transform is necessarily unique despite collinearity.
\item \(I_t\) is geometrically degenerate and does not uniquely constrain a general 2D affine transform.
\item An all-inlier sample is always sufficient; degeneracy can occur only when at least one correspondence is wrong.
\item Repeated least-squares refitting of the same three pairs automatically restores the missing rank.
\end{enumerate}
\end{minipage}
\par\vspace{12pt}
\begin{minipage}{\linewidth}
\textbf{Q16 — Multiple choice}\par\smallskip
Increasing the RANSAC inlier threshold gives a model many more counted inliers. Why is this not automatically evidence of a better geometric model?
\begin{enumerate}[label=\textbf{\Alph*.}]
\item A larger threshold always improves localization accuracy.
\item A threshold has no influence on consensus size.
\item The larger threshold can admit poor matches and outliers as inliers.
\item Every counted inlier has exactly zero residual by definition.
\end{enumerate}
\end{minipage}
\par\vspace{12pt}
\begin{minipage}{\linewidth}
\textbf{Q17 — Multiple choice}\par\smallskip
A stationary scene undergoes a moving cast shadow. A brightness-constancy flow method reports motion in the shadow region. Which interpretation is correct?
\begin{enumerate}[label=\textbf{\Alph*.}]
\item The appearance-constancy assumption for a physical surface point is violated; apparent image motion need not equal physical surface motion.
\item Increasing the window size always establishes that reported image motion corresponds to physical surface motion.
\item The reported vector must be a sign error,\allowbreak{} because a static scene guarantees zero estimated flow regardless of illumination.
\item Every measured optical flow vector necessarily equals a 3D object's velocity.
\end{enumerate}
\end{minipage}
\par\vspace{12pt}
\begin{minipage}{\linewidth}
\textbf{Q18 — Multiple choice}\par\smallskip
Why can a coarse-to-fine flow method still fail on a tiny fast-moving object?
\begin{enumerate}[label=\textbf{\Alph*.}]
\item Small objects always have zero displacement at full resolution.
\item The coarse-level flow already gives exact depth in metres.
\item Pyramids make every flow observable regardless of texture or visibility.
\item The object can disappear or become indistinct at coarse levels,\allowbreak{} so its motion is not initialized correctly.
\end{enumerate}
\end{minipage}
\par\vspace{12pt}
\begin{minipage}{\linewidth}
\textbf{Q19 — Multiple choice}\par\smallskip
An online tracker updates its appearance model from its own predicted box. What is a major risk during occlusion?
\begin{enumerate}[label=\textbf{\Alph*.}]
\item Any online update guarantees improved test accuracy.
\item Incorrect predictions can contaminate the model and cause tracking drift.
\item Occlusion automatically provides perfect new positive labels.
\item The model can never change because the initial template is fixed by definition.
\end{enumerate}
\end{minipage}
\par\vspace{12pt}
\begin{minipage}{\linewidth}
\textbf{Q20 — Multiple choice}\par\smallskip
A camera translates while observing a scene with nearby and distant objects. One fitted panorama homography aligns the distant background but not the foreground. What is the main cause?
\begin{enumerate}[label=\textbf{\Alph*.}]
\item Depth-dependent parallax; one global homography need not align the nonplanar scene.
\item The foreground and background require opposite matrix multiplication orders,\allowbreak{} even when the transform direction is fixed.
\item The blend ramp is too narrow; widening it guarantees correction of every geometric mismatch.
\item H has the wrong homogeneous scale; multiplying all its entries by a different nonzero scalar corrects foreground without moving background.
\end{enumerate}
\end{minipage}
\par\vspace{12pt}
\begin{minipage}{\linewidth}
\textbf{Q21 — Multiple choice}\par\smallskip
H maps source image coordinates to a reference image. A translation T then moves reference coordinates into a larger panorama canvas. With column vectors,\allowbreak{} which matrix maps source directly to the canvas?
\begin{enumerate}[label=\textbf{\Alph*.}]
\item \(TH\)
\item \(HT\)
\item \(H+T\)
\item \(H^{\mathsf T}T\),\allowbreak{} regardless of the coordinate convention
\end{enumerate}
\end{minipage}
\par\vspace{12pt}
\begin{minipage}{\linewidth}
\textbf{Q22 — Multiple choice}\par\smallskip
Which classifier and prediction does this code implement for binary class labels? Squared distances preserve Euclidean neighbour order. Provided routine: \texttt{np.argsort(d2)} returns indices sorted by increasing d2.


\begin{Code}
import numpy as np
X = np.array([[0,0], [0,1], [1,0], [3,3]], dtype=float)
y = np.array([1,0,0,1])
query = np.array([.1,.1])
d2 = ((X-query)**2).sum(axis=1)
nearest = np.argsort(d2)[:3]
prediction = int(y[nearest].sum() > 1)
\end{Code}

\begin{enumerate}[label=\textbf{\Alph*.}]
\item K-means predicts class 0 after moving three learned centres; the array y is used only as an initialization seed.
\item Unweighted 3-nearest-neighbour voting predicts class 1 because the closest example's label outweighs both others.
\item Unweighted 3-nearest-neighbour majority voting predicts class 0,\allowbreak{} although the closest example has class 1.
\item 1-nearest-neighbour classification predicts class 1; the two other selected labels never contribute.
\end{enumerate}
\end{minipage}
\par\vspace{12pt}
\begin{minipage}{\linewidth}
\textbf{Q23 — Multiple choice}\par\smallskip
For pixel clustering,\allowbreak{} why might features \([R,G,B,\alpha x,\alpha y]\) give different segments from \([R,G,B]\)?
\begin{enumerate}[label=\textbf{\Alph*.}]
\item K-means always returns connected regions even with colour-only features.
\item Adding x,\allowbreak{}y guarantees the true semantic labels without supervision.
\item Spatial terms encourage location coherence,\allowbreak{} and \(\alpha\) sets their relative importance against colour differences.
\item Feature scaling cannot affect Euclidean distances.
\end{enumerate}
\end{minipage}
\par\vspace{12pt}
\begin{minipage}{\linewidth}
\textbf{Q24 — Multiple choice}\par\smallskip
A training set contains label noise and is not linearly separable. Why use a soft-margin rather than a hard-margin linear SVM?
\begin{enumerate}[label=\textbf{\Alph*.}]
\item \(I_t\) removes every mislabeled point with certainty.
\item \(I_t\) guarantees infinite margin for nonseparable data.
\item \(I_t\) permits margin violations while balancing fit against regularization.
\item \(I_t\) ignores the weight norm entirely and only memorizes pixels.
\end{enumerate}
\end{minipage}
\par\vspace{12pt}
\begin{minipage}{\linewidth}
\textbf{Q25 — Multiple choice}\par\smallskip
A CNN outputs softmax class scores that sum to one. Which conclusion is NOT guaranteed merely by applying softmax?
\begin{enumerate}[label=\textbf{\Alph*.}]
\item Increasing one logit relative to the others increases its softmax probability.
\item Each output is nonnegative.
\item The outputs sum to one,\allowbreak{} up to numerical precision.
\item The probabilities are well calibrated to actual correctness frequencies.
\end{enumerate}
\end{minipage}
\par\vspace{12pt}
\begin{minipage}{\linewidth}
\textbf{Q26 — Multiple choice}\par\smallskip
Consider this training sequence from the kind of loop taught in Lab5. What is the purpose of the first line? No executable PyTorch code is required.


\begin{Code}
optimizer.zero_grad()
loss.backward()
optimizer.step()
\end{Code}

\begin{enumerate}[label=\textbf{\Alph*.}]
\item Disable gradient recording for this batch while retaining earlier batches' stored gradients.
\item Reset learned model parameters to zero,\allowbreak{} ensuring that each batch starts from the same initialized network.
\item Normalize every stored gradient to unit norm,\allowbreak{} so the learning rate becomes irrelevant to the next update.
\item Prevent unintended accumulation of earlier batches' gradients before this batch's update.
\end{enumerate}
\end{minipage}
\par\vspace{12pt}
\begin{minipage}{\linewidth}
\textbf{Q27 — Multiple choice}\par\smallskip
A residual branch changes feature shape from \(H\times W\times C\) to \(\frac{H}{2}\times\frac{W}{2}\times2C\). What is required for an additive shortcut?
\begin{enumerate}[label=\textbf{\Alph*.}]
\item A compatible shortcut,\allowbreak{} such as a stride-2 projection,\allowbreak{} so the tensors have matching shapes.
\item Remove all nonlinearities from the entire network.
\item Add the mismatched tensors with no mapping; addition always changes the channel count automatically.
\item Concatenate them and call the result an identity addition.
\end{enumerate}
\end{minipage}
\par\vspace{12pt}
\begin{minipage}{\linewidth}
\textbf{Q28 — Multiple choice}\par\smallskip
Multiply both a predicted box and its ground-truth box coordinates by the same positive scale. Which property is correct?
\begin{enumerate}[label=\textbf{\Alph*.}]
\item Both losses must be identical for every pair of boxes.
\item IoU doubles whenever object area doubles.
\item Coordinate error becomes scale invariant just by summing its four components.
\item IoU is unchanged,\allowbreak{} while an unnormalized squared coordinate-error loss generally scales by the square of that factor.
\end{enumerate}
\end{minipage}
\par\vspace{12pt}
\begin{minipage}{\linewidth}
\textbf{Q29 — Multiple choice}\par\smallskip
For an ideal rectified stereo pair with parallel optical axes and a horizontal baseline,\allowbreak{} why are epipolar lines parallel?
\begin{enumerate}[label=\textbf{\Alph*.}]
\item The epipoles are at infinity in the horizontal baseline direction.
\item The fundamental matrix maps every point to the image's principal point.
\item Both cameras have exactly the same optical centre.
\item Every scene point has the same depth.
\end{enumerate}
\end{minipage}
\par\vspace{12pt}
\begin{minipage}{\linewidth}
\textbf{Q30 — Multiple choice}\par\smallskip
Which description of a standard multi-view reconstruction pipeline is correct?
\begin{enumerate}[label=\textbf{\Alph*.}]
\item Bundle adjustment modifies only colour histograms and never camera parameters.
\item SfM estimates camera poses and sparse structure; bundle adjustment refines poses and points jointly; MVS can densify using estimated cameras.
\item SfM requires only one camera image with no priors or known geometry.
\item MVS always estimates metric scale from arbitrary uncalibrated image content alone.
\end{enumerate}
\end{minipage}
\par\vspace{12pt}
\clearpage\section*{Part B — Numerical values}
\begin{minipage}{\linewidth}
\textbf{Q31 — Numerical}\par\smallskip
An original camera has \(f_x=800,\;f_y=600,\;c_x=320,\;c_y=240\) and zero skew. Crop away 40 pixels from the left and 20 from the top; then halve both spatial dimensions. Use continuous coordinates with no half-pixel offsets. A camera-frame point is \(\begin{bmatrix}0.5&1&2\end{bmatrix}\). Report its final image v coordinate,\allowbreak{} using the virtual image plane in front of the camera.
\par\medskip\textbf{Answer:} \rule{55mm}{0.4pt}
\end{minipage}
\par\vspace{12pt}
\begin{minipage}{\linewidth}
\textbf{Q32 — Numerical}\par\smallskip
For the RGB image below,\allowbreak{} compute normalized RGB PER PIXEL,\allowbreak{} then report the spatial mean of its normalized red channel. For black pixels,\allowbreak{} define all normalized channels as zero.


\begin{Code}
import numpy as np
I = np.array([[[6,1,3], [2,6,2]],
              [[0,0,0], [1,1,2]]], dtype=float)
s = I.sum(axis=2, keepdims=True)
safe = np.where(s == 0, 1, s)
normalized = I / safe
answer = normalized[:, :, 0].mean()
\end{Code}


Provided function: \texttt{np.where(condition,\allowbreak{} a,\allowbreak{} b)} selects a where the condition is true and b elsewhere. Black pixels are included in the spatial mean.
\par\medskip\textbf{Answer:} \rule{55mm}{0.4pt}
\end{minipage}
\par\vspace{12pt}
\begin{minipage}{\linewidth}
\textbf{Q33 — Numerical}\par\smallskip
Let estimated height gradients be \(p(x,y)=\frac{\partial z}{\partial x}=y\) and \(q(x,y)=\frac{\partial z}{\partial y}=2x\). Start with \(z(0,0)=0\). Integrate continuously to (2,\allowbreak{}3) along path A: x first then y; and path B: y first then x. Report the absolute difference between the two resulting heights.
\par\medskip\textbf{Answer:} \rule{55mm}{0.4pt}
\end{minipage}
\par\vspace{12pt}
\begin{minipage}{\linewidth}
\textbf{Q34 — Numerical}\par\smallskip
Use calibrated Lambertian intensity \(I=\rho\max(n\cdot l,0)\),\allowbreak{} with unit vectors \(n=\begin{bmatrix}2/3&-1/3&2/3\end{bmatrix}^{\mathsf T}\) and \(l=\begin{bmatrix}0&0.6&0.8\end{bmatrix}^{\mathsf T}\),\allowbreak{} \(\rho=0.75\),\allowbreak{} and unit light strength/\allowbreak{}gain. Report I.
\par\medskip\textbf{Answer:} \rule{55mm}{0.4pt}
\end{minipage}
\par\vspace{12pt}
\begin{minipage}{\linewidth}
\textbf{Q35 — Numerical}\par\smallskip
A 9 by 9 binary image is zero except for a 3 by 3 all-ones block at rows 3..5 and columns 3..5 (inclusive,\allowbreak{} zero-based),\allowbreak{} and a single foreground pixel at (0,\allowbreak{}0). Apply opening with a centred 3 by 3 all-ones element,\allowbreak{} zero outside the image,\allowbreak{} retaining the same shape after each stage. Report the final foreground-pixel count.
\par\medskip\textbf{Answer:} \rule{55mm}{0.4pt}
\end{minipage}
\par\vspace{12pt}
\begin{minipage}{\linewidth}
\textbf{Q36 — Numerical}\par\smallskip
Use the median of three entries: \texttt{med(v)=\allowbreak{}sorted(v)[1]}. Let \(a=\begin{bmatrix}0&0&10\end{bmatrix}\) and \(b=\begin{bmatrix}0&10&0\end{bmatrix}\),\allowbreak{} with elementwise addition. Report \(\operatorname{med}(a+b)-\operatorname{med}(a)-\operatorname{med}(b)\) as an exact integer.
\par\medskip\textbf{Answer:} \rule{55mm}{0.4pt}
\end{minipage}
\par\vspace{12pt}
\begin{minipage}{\linewidth}
\textbf{Q37 — Numerical}\par\smallskip
For \(f=\begin{bmatrix}1&0&2&0\end{bmatrix}\) and \(h=\begin{bmatrix}0&1&0&1\end{bmatrix}\),\allowbreak{} use 4-point CIRCULAR convolution \(c_k=\sum_{n=0}^{3}f_n h_{(k-n)\bmod4}\). Report \(c_3\) as an exact integer. This is the result of multiplying their unpadded 4-point DFTs and applying the inverse DFT.
\par\medskip\textbf{Answer:} \rule{55mm}{0.4pt}
\end{minipage}
\par\vspace{12pt}
\begin{minipage}{\linewidth}
\textbf{Q38 — Numerical}\par\smallskip
The grid below contains gradient magnitudes AFTER non-maximum suppression. Use 8-neighbour Canny hysteresis: strong means \(\ge9\); weak means \(\ge4\text{ and }<9\); lower values are rejected. Retain every weak pixel connected through weak/\allowbreak{}strong pixels to a strong pixel. Report the retained pixel count.


\[
\begin{bmatrix}10&6&0&0&5&5&0\\0&6&0&0&0&5&0\\0&6&6&0&0&0&0\end{bmatrix}
\]

\par\medskip\textbf{Answer:} \rule{55mm}{0.4pt}
\end{minipage}
\par\vspace{12pt}
\begin{minipage}{\linewidth}
\textbf{Q39 — Numerical}\par\smallskip
Use equal weights for a four-pixel window with \(I_x=\begin{bmatrix}1&0&1&2\end{bmatrix}\) and \(I_y=\begin{bmatrix}0&2&1&0\end{bmatrix}\). Construct
\[M=\begin{bmatrix}\sum I_x^2&\sum I_x I_y\\\sum I_x I_y&\sum I_y^2\end{bmatrix}.\]
 Report the Shi–Tomasi score,\allowbreak{} the smaller eigenvalue. Provided formula for symmetric
\[\begin{bmatrix}a&b\\b&d\end{bmatrix}:\]
 \(\lambda_{\min}=\frac{a+d-\sqrt{(a-d)^2+4b^2}}{2}\).
\par\medskip\textbf{Answer:} \rule{55mm}{0.4pt}
\end{minipage}
\par\vspace{12pt}
\begin{minipage}{\linewidth}
\textbf{Q40 — Numerical}\par\smallskip
A flow at one-eighth spatial resolution is \([u,v]=\begin{bmatrix}1.75&-0.5\end{bmatrix}\) in COARSE pixels. Upsample to full resolution with the correct vector-unit scaling,\allowbreak{} then add a full-resolution residual flow \(\begin{bmatrix}0.25&0.75\end{bmatrix}\). Report the Euclidean magnitude of the resulting full-resolution vector.
\par\medskip\textbf{Answer:} \rule{55mm}{0.4pt}
\end{minipage}
\par\vspace{12pt}
\begin{minipage}{\linewidth}
\textbf{Q41 — Numerical}\par\smallskip
An affine RANSAC uses minimal samples of 3 point pairs and independently sampled inlier fraction \(w=0.4\). Assume all-inlier samples are nondegenerate. Use \(N=\left\lceil\frac{\log(1-p)}{\log(1-w^3)}\right\rceil\). Report the minimum integer N for \(p=0.999\).
\par\medskip\textbf{Answer:} \rule{55mm}{0.4pt}
\end{minipage}
\par\vspace{12pt}
\begin{minipage}{\linewidth}
\textbf{Q42 — Numerical}\par\smallskip
A nonsingular homography maps these four source points to targets:


\[
\begin{aligned}(0,0)&\mapsto (1,2)\\(1,0)&\mapsto (3,2)\\(0,1)&\mapsto (1,4)\\(1,1)&\mapsto (3,3)\end{aligned}
\]


Use
\[H=\begin{bmatrix}a&b&c\\d&e&f\\g&h&1\end{bmatrix},\]
 so \(x'=\frac{ax+by+c}{gx+hy+1},\quad y'=\frac{dx+ey+f}{gx+hy+1}\). Rearrange the eight scalar equations into a linear system. You may use the provided least-squares routine \texttt{np.linalg.lstsq(A,\allowbreak{} rhs,\allowbreak{} rcond=\allowbreak{}None)[0]} to solve a full-rank system. Report the target x coordinate for source (0.5,\allowbreak{}0.5).
\par\medskip\textbf{Answer:} \rule{55mm}{0.4pt}
\end{minipage}
\par\vspace{12pt}
\begin{minipage}{\linewidth}
\textbf{Q43 — Numerical}\par\smallskip
A matcher returns SQUARED Euclidean distances \([D_1,D_2]\),\allowbreak{} nearest first. Apply Lowe's test to the ordinary distances: \(\frac{\sqrt{D_1}}{\sqrt{D_2}}<0.7\). For the rows below,\allowbreak{} report the largest ordinary nearest distance among accepted matches.


\[
\begin{bmatrix}4&9\\9&16\\1&4\\25&36\\0.25&1\end{bmatrix}
\]

\par\medskip\textbf{Answer:} \rule{55mm}{0.4pt}
\end{minipage}
\par\vspace{12pt}
\begin{minipage}{\linewidth}
\textbf{Q44 — Numerical}\par\smallskip
A 1D visual vocabulary has centres \(\begin{bmatrix}0&4&8\end{bmatrix}\). Quantize descriptors \(\begin{bmatrix}1&3&5&7&8&9\end{bmatrix}\) to their nearest centre using Euclidean distance. Build counts,\allowbreak{} then L1-normalize by the number of descriptors. Report the normalized bin value for centre 8.
\par\medskip\textbf{Answer:} \rule{55mm}{0.4pt}
\end{minipage}
\par\vspace{12pt}
\begin{minipage}{\linewidth}
\textbf{Q45 — Numerical}\par\smallskip
A binary classifier predicts positive iff \(\text{score}\ge0.6\). True labels are \(\begin{bmatrix}1&0&1&0&1&0\end{bmatrix}\) and scores are \(\begin{bmatrix}0.9&0.8&0.7&0.6&0.5&0.4\end{bmatrix}\),\allowbreak{} in the same order. Report F1 for the positive class as a fraction.
\par\medskip\textbf{Answer:} \rule{55mm}{0.4pt}
\end{minipage}
\par\vspace{12pt}
\begin{minipage}{\linewidth}
\textbf{Q46 — Numerical}\par\smallskip
One class has 4 ground-truth objects. After score-ordered,\allowbreak{} one-to-one matching,\allowbreak{} ranked detections are [FP,\allowbreak{}TP,\allowbreak{}FP,\allowbreak{}TP,\allowbreak{}TP]. Use ALL-POINTS INTERPOLATED AP: sum each positive recall increment times the maximum precision at that recall or any later recall. Precision beyond maximum achieved recall is zero. Report AP as a fraction. This is NOT 11-point AP or raw uninterpolated AP.
\par\medskip\textbf{Answer:} \rule{55mm}{0.4pt}
\end{minipage}
\par\vspace{12pt}
\begin{minipage}{\linewidth}
\textbf{Q47 — Numerical}\par\smallskip
A residual block has an identity shortcut and this dense bottleneck branch: \(1\times1\) convolution \(64\to16\) channels; \(3\times3\) convolution \(16\to16\); \(1\times1\) convolution \(16\to64\). Include one bias per output channel in EACH convolution. There is no normalization layer or trainable shortcut. Report total trainable parameters.
\par\medskip\textbf{Answer:} \rule{55mm}{0.4pt}
\end{minipage}
\par\vspace{12pt}
\begin{minipage}{\linewidth}
\textbf{Q48 — Numerical}\par\smallskip
A standard 3D convolution processes a video tensor with 2 input channels and uses 5 output channels. Each filter has temporal x height x width size \(3\times3\times3\). Include one bias per output channel. No groups. Report the trainable parameter count; the video length and spatial image size do not enter this count.
\par\medskip\textbf{Answer:} \rule{55mm}{0.4pt}
\end{minipage}
\par\vspace{12pt}
\begin{minipage}{\linewidth}
\textbf{Q49 — Numerical}\par\smallskip
Use the ORIGINAL YOLO v1 confidence convention: \(\text{box confidence}=\Pr(\text{object})\cdot\operatorname{IoU}(\text{predicted box},\mathrm{GT})\); \(\text{class-specific score}=\Pr(\text{class}\mid\text{object})\cdot\text{box confidence}\). Given \(\Pr(\text{object})=0.75\),\allowbreak{} \(\mathrm{IoU}=0.6\) and \(\Pr(\text{class}\mid\text{object})=0.8\),\allowbreak{} report the class-specific score.
\par\medskip\textbf{Answer:} \rule{55mm}{0.4pt}
\end{minipage}
\par\vspace{12pt}
\begin{minipage}{\linewidth}
\textbf{Q50 — Numerical}\par\smallskip
For rectified stereo,\allowbreak{} \(Z=\frac{fB}{d}\). Let \(f=900\) pixels,\allowbreak{} \(B=0.2\) metres,\allowbreak{} \(d=12\) pixels,\allowbreak{} and small disparity-error magnitude \(\delta_d=0.25\) pixel. Use FIRST-ORDER error propagation \(\lvert\delta_Z\rvert\approx\frac{fB}{d^2}\lvert\delta_d\rvert\). Report the estimated depth-error magnitude in metres,\allowbreak{} not the exact finite-difference error.
\par\medskip\textbf{Answer:} \rule{55mm}{0.4pt}
\end{minipage}
\par\vspace{12pt}

\end{document}
